\begin{titlepage}
\begin{tikzpicture}[overlay, remember picture, shift={(current page.center)}, yshift=0.1\paperheight]
    \def\Ox{0}
    \def\Oy{0}
    
    \draw[blue!30, line width=0.3pt] (\Ox,\Oy) circle (8cm);
    \draw[blue!20, line width=0.3pt] (\Ox,\Oy) circle (6.5cm);
    \draw[blue!10, line width=0.3pt] (\Ox,\Oy) circle (5cm);
    
    \def\R{8}
    \def\omega{30}
    \pgfmathsetmacro{\theta}{\omega}
    \pgfmathsetmacro{\Px}{\R*cos(\theta)}
    \pgfmathsetmacro{\Py}{\R*sin(\theta)}
    
    \filldraw[red!60] (\Px,\Py) circle (2.5pt);
    
    \def\r{1.8}
    \def\omegai{120}
    \pgfmathsetmacro{\thetai}{\omegai}
    \pgfmathsetmacro{\Pix}{\Px + \r*cos(\thetai)}
    \pgfmathsetmacro{\Piy}{\Py + \r*sin(\thetai)}
    
    \draw[red!50, line width=0.4pt] (\Px,\Py) circle (\r);
    \filldraw[red!50] (\Pix,\Piy) circle (2.0pt);
    
    \def\rr{0.8}
    \def\omegaii{210}
    \pgfmathsetmacro{\thetaii}{\omegaii}
    \pgfmathsetmacro{\Piix}{\Pix + \rr*cos(\thetaii)}
    \pgfmathsetmacro{\Piiy}{\Piy + \rr*sin(\thetaii)}
    
    \draw[orange!70, line width=0.4pt] (\Pix,\Piy) circle (\rr);
    \filldraw[orange!80] (\Piix,\Piiy) circle (1.6pt);
    
    \def\rrr{0.35}
    \def\omegaiii{330}
    \pgfmathsetmacro{\thetaiii}{\omegaiii}
    \pgfmathsetmacro{\Piiix}{\Piix + \rrr*cos(\thetaiii)}
    \pgfmathsetmacro{\Piiiy}{\Piiy + \rrr*sin(\thetaiii)}
    
    \draw[green!60!black, line width=0.4pt] (\Piix,\Piiy) circle (\rrr);
    \filldraw[green!70!black] (\Piiix,\Piiiy) circle (1.2pt);
    
    \draw[blue!15, line width=0.2pt, densely dotted] (\Ox,\Oy) -- (\Px,\Py);
    \draw[red!15, line width=0.2pt, densely dotted] (\Px,\Py) -- (\Pix,\Piy);
    \draw[orange!15, line width=0.2pt, densely dotted] (\Pix,\Piy) -- (\Piix,\Piiy);
    \draw[green!15, line width=0.2pt, densely dotted] (\Piix,\Piiy) -- (\Piiix,\Piiiy);
    
    \foreach \k in {0,1,...,59}
    {
        \filldraw[blue!40] ({360/60*\k}:8cm) circle (1.2pt);
        \filldraw[blue!30] ({360/60*\k + 3}:6.5cm) circle (1.0pt);
        \filldraw[blue!20] ({360/60*\k + 6}:5cm) circle (0.8pt);
    }

    \foreach \k in {0,1,...,11}
    {
        \filldraw[red!50] ({360/12*\k}:9cm) circle (2.0pt);
    }

    \foreach \k in {0,1,...,7}
    {
        \filldraw[green!50!black] ({360/8*\k + 11.25}:7.2cm) circle (1.8pt);
    }

    \foreach \k in {0,1,...,4}
    {
        \filldraw[orange!60] ({360/5*\k + 18}:5.8cm) circle (1.6pt);
    }
\end{tikzpicture}

\begin{tikzpicture}[overlay, remember picture, shift={(current page.south)}, yshift=0.5cm]
    \node[font=\footnotesize, text=black!30] at (0,0) 
        [yshift=2.0cm] {$\mathcal{P}_T = \{(A_k, \omega_k, \phi_k) \mid A_k \in \mathbb{R}^+, \omega_k = 2\pi k/T, \phi_k \in [0, 2\pi), k = 0, 1, \dots, T-1\}$};
    
    \node[font=\footnotesize, text=black!30] at (0,0) 
        [yshift=1.5cm] {$\mathcal{K}^{(d)}\{f, g\} = \sum_{n}\sum_{m} A_{n,m} \cdot \operatorname{Sim}^{(d)}(\theta_n, \theta_m) \cdot e^{i(\Omega_{n,m} x + \Phi_{n,m})}$};
\end{tikzpicture}

\vspace*{0.5cm}

\begin{center}

{\fontsize{36pt}{42pt}\selectfont \textbf{Periodic Mathematics}}

\vspace{0.1cm}

\begin{center}
\itshape
``The tendency that could guide could not always guide. The symbol that could name could not always name.'' \\[0.2cm]
``The tendency evolves into the one, the one evolves into the two, the two evolves into the three, the three evolves into the everything.'' \\[0.2cm]
— Tao Te Ching
\end{center}

\vspace{0.1cm}

\begin{center}
\itshape
``Observe the principle of the Universe, follow the course of the Universe. That is all.'' \\[0.2cm]
— Yin Fu Ching
\end{center}

\vspace{0.1cm}

\begin{center}
\itshape
``Mathematics is the language with which God has written the universe.'' \\[0.2cm]
— Galileo Galilei
\end{center}

\vspace{0.1cm}

\begin{center}
\itshape
``This discovery has convinced me that Fuxi was not only the founder of Chinese philosophy, but also one of the most ancient sages of all mankind. The system of symbols he left behind contains a universal writing---a character capable of expressing every thought.'' \\[0.2cm]
``It seems that the Chinese, in remote antiquity, already possessed a universal writing, of which the Eight Trigrams are the vestige. It is only that later men forgot its true meaning, and used it merely for divination.'' \\[0.2cm]
— Gottfried Wilhelm Leibniz
\end{center}

\vspace{0.1cm}

\begin{center}
\itshape
``Kindred sounds resonate; kindred spirits gravitate.'' \\[0.2cm]
``Nothing flat does not slope; nothing gone does not return.'' \\[0.2cm]
— I Ching
\end{center}

\vspace{0.1cm}

\begin{center}
\itshape
``The void evolves into the cosmos which evolves into the tendency that is limited in bounds. The clarity evaporates into the sky while the mud sinks into the ground. The clarity evolves into the positive while the mud evolves into the negative. The positive and the negative evolves into the seasons and the seasons evolves into the myriad things.'' \\[0.2cm]
— Huainanzi (Treatise on the Patterns of Heaven)
\end{center}

\vspace{0.1cm}

\begin{center}
``In birth, death is already present; in death, birth is already present. In what is deemed possible, the impossible is already present; in what is deemed impossible, the possible is already present. Right depends on wrong, and wrong depends on right.''
— Zhuangzi, (On the Equality of Things)
\end{center}

\begin{center}
\itshape
``The great similarity that differs with small similarity is called small similarity-and-difference. Similarity and difference among the everything called great similarity-and-difference.'' \\[0.2cm]
— Zhuangzi (All Under Heaven)
\end{center}

\vfill

\end{center}
\end{titlepage}